\ifdefined\CRevisionRound\else\def\CRevisionRound{2}\fi
\documentclass[10pt]{article}
\pdfvariable suppressoptionalinfo 767
\usepackage[a4paper,margin=20mm]{geometry}
\usepackage{amsmath,amssymb,amsthm,booktabs}
\usepackage{fontspec,unicode-math}
\setmainfont{TeX Gyre Pagella}
\setmathfont{TeX Gyre Pagella Math}
\setmonofont{Latin Modern Mono}
\newfontfamily\cjkfont{Droid Sans Fallback}
\usepackage[hidelinks]{hyperref}
\newcommand{\ii}{\mathrm i}
\newcommand{\e}{\mathrm e}
\newcommand{\dd}{\,\mathrm d}
\newcommand{\R}{\mathbb R}
\newcommand{\C}{\mathbb C}
\newcommand{\Z}{\mathbb Z}
\DeclareMathOperator{\Ree}{Re}
\DeclareMathOperator{\dom}{Dom}
\newtheorem{theorem}{Theorem}
\newtheorem{proposition}{Proposition}
\ifcase\CRevisionRound
\title{An Impedance String in Physical Time:\\Complete Spectrum and Transparent Extinction}
\or
\title{An Impedance String in Physical Time:\\Empty Spectrum and Exact Transparent Pseudospectra}
\else
\title{An Impedance String in Physical Time:\\Empty Spectrum, Exact Pseudospectra and Determinant Boundaries}
\fi
\author{HCS-C396 theorem and reproducibility package}
\date{5 September 2026}
\begin{document}
\emergencystretch=2em
\maketitle
\begin{abstract}
\ifcase\CRevisionRound
We reconstruct the full energy-space dynamics of a homogeneous string with
one fixed endpoint and one finite nonnegative impedance. A unitary
characteristic unfolding preserves physical time and transforms the wave
generator into a derivative with one scalar boundary condition.
For every nonzero signed round-trip coefficient, the entire spectrum is a
simple bilateral arithmetic ladder with a complete Riesz basis.
The semigroup norm is an exact staircase, including all round-trip endpoints.
At transparent impedance the generator spectrum is empty, yet the semigroup
has norm one until precisely twice the length divided by the wave speed;
after that time it vanishes. An explicit all-plane resolvent verifies that
empty spectrum is an operator theorem, not a failed eigenvalue search.
Reciprocal positive impedances give the same norm and different spectral
ladders. The classical absorbing-wave mechanism retains its literature
ownership. This draft closes the domain, clock, complete spectrum and all-time
law, while a finite exact and high-precision ledger audits their constants.
\or
We give an exact transparent resolvent analysis for the homogeneous
impedance string, retaining its complete energy-space domain and physical
clock. The generator has empty spectrum at perfect absorption, although its
semigroup remains norm one up to the sharp round-trip extinction time.
The resolvent norm is reduced to the lowest eigenvalue of a separated
Sturm--Liouville problem with a parameter-dependent Robin condition.
Positive interior eigenfunctions select the correct trigonometric,
critical and hyperbolic branches without relying on sampled singular values.
The resulting norm decreases strictly with the real part of the spectral
parameter, so every transparent pseudospectrum is an exactly specified left
half-plane. At zero spectral parameter its norm is twice the round-trip time
divided by pi. Independent Green-kernel, Rayleigh quotient and complex gauge
integrals test the formulas. These finite checks do not replace the
infinite-dimensional proof, and no new ownership of the classical
absorbing-wave spectral mechanism is asserted.
\else
We assemble a domain-complete theorem for a finite homogeneous string with
a fixed endpoint and a nonnegative finite impedance at the other endpoint.
Exact unitary unfolding gives the complete nontransparent spectrum, a Riesz
basis and the physical-time semigroup norm. At transparency the spectrum
is empty, but the evolution remains norm one until the sharp round-trip
extinction time. A separated singular-value problem gives every transparent
pseudospectrum through a three-branch scalar formula. We then distinguish
compact resolvent from trace-class resolvent: the transparent inverse has
half-integer singular values, is Hilbert--Schmidt but not trace class, and
is quasinilpotent. Its explicitly named order-two regularized determinant is
identically one. Before extinction the semigroup is noncompact, so it has
no ordinary trace. Reciprocal impedances further separate decay norm from
spectral location. The classical formulas are attributed; independent exact
and numerical lanes, strict typed gates and repaired-hash attacks delimit
the package's reproducible contribution. No target arithmetic or target
self-adjoint operator is inferred from the native wave generator.
\fi
\end{abstract}
\noindent\textbf{Keywords:} impedance string; physical-time semigroup;
empty spectrum; exact pseudospectrum; Volterra resolvent; determinant boundary.

\medskip
\noindent{\cjkfont\small\raggedright 中文摘要：\par
\ifcase\CRevisionRound
本文完整整理一端固定、另一端具有有限非负阻抗的均匀弦动力学。
精确等距展开保留物理时间，将波动生成元化为带单个边界条件的一阶导数。
非零往返反射系数给出全部简单谱点及完备基，半群范数呈精确阶梯。
透明边界下生成元没有谱点，但演化在完整往返时间以前仍具有单位范数，
此后才完全消亡。全文明确区分无限维证明与有限回归检查，不宣称经典公式的首创。
\or
本文在完整物理域与时间规范下，推导透明阻抗弦的精确预解式范数。
相应奇异值问题具有参数相关边界条件，最低正特征函数确定三角、临界及双曲三种分支。
范数沿谱参数实部严格下降，因此每个透明拟谱均为精确定义的左半平面。
独立积分检查核对作用公式与归一化，有限计算不替代全谱及全时间定理。
\else
本文给出阻抗弦从全谱到透明空谱的完整定理，并明确算子理想与行列式的边界。
透明预解式是紧的平方可和算子，但不是迹类，其二阶正则化行列式恒等于一。
半群在有限时消亡前不紧，因此不存在该阶段的普通算子迹。
互为倒数的正阻抗具有相同衰减范数，却具有不同谱列。
经典机制保留原文所有权，严格核验仅支撑本包的复算边界，不产生目标算术或目标自伴算子。
\fi
\par 关键词：阻抗弦；物理时间；空谱；精确拟谱；预解式；行列式边界。
\par}

\section{One source, three different spectral questions}
Perfect absorption separates three quantities that are often conflated:
the spectrum of an unbounded generator, the norm of its evolution, and the
norm of its resolvent. Our source is the homogeneous wave equation
\begin{equation}
 u_{tt}=c^2u_{xx},\quad 0<x<L,\qquad
 u(0,t)=0,\quad u_x(L,t)+(\eta/c)u_t(L,t)=0,
 \label{eq:pde}
\end{equation}
where $L,c>0$ and $0\leq\eta<\infty$. The time in this equation is never
replaced by a spectral or target-dependent clock. Define
\begin{equation}
 \ell=L/c,\qquad \tau=2L/c,\qquad q=(\eta-1)/(\eta+1)\in[-1,1).
 \label{eq:q}
\end{equation}
The signed coefficient $q$ includes the reflection at the fixed endpoint;
it is the negative of the right-end reflection coefficient alone.

Driscoll and Trefethen~\cite{driscoll1996} established the classical
absorbing-wave mechanism, the exact decay staircase and the empty-spectrum
transparent case. Their Theorem~1 gives resolvent bounds and invariances.
We retain that ownership. The package contribution is a self-contained
domain-and-clock reconstruction joined to independently executable evidence
and explicit operator-category limitations, not a claim of literature priority.
The source theorem is substantially different from free-wave observability:
there is no boundary control input, observation inequality or control cost.

\section{The physical domain and its unitary unfolding}
Put $v=u_t$ and $p=cu_x$. On the complex Hilbert space
\begin{equation}
 \mathcal H=\{u\in H^1(0,L):u(0)=0\}\times L^2(0,L),\qquad
 \|(u,v)\|_E^2=\int_0^L(|p|^2+|v|^2)\dd x,
\end{equation}
the generator is $G_\eta(u,v)=(v,c^2u_{xx})$. Its domain consists exactly
of $u\in H^2$, $v\in H^1$, with $u(0)=v(0)=0$ and the right boundary
in~\eqref{eq:pde}. Physical energy is one half of the squared norm.

\begin{theorem}[Exact domain equivalence]
The following map is unitary from $\mathcal H$ to $L^2(0,\tau)$:
\begin{equation}
 U(u,v)(s)=\sqrt{c/2}
 \begin{cases}
 (p-v)(L-cs),&0<s<\ell,\\
 (p+v)(c(s-\ell)),&\ell<s<\tau.
 \end{cases}
 \label{eq:unfold}
\end{equation}
It sends the entire physical generator, including its domain, to
\begin{equation}
 A_qw=w',\qquad \dom(A_q)=\{w\in H^1(0,\tau):w(\tau)=qw(0)\}.
 \label{eq:domain}
\end{equation}
\end{theorem}
\begin{proof}
The two changes of variables in~\eqref{eq:unfold}, and
$|p-v|^2+|p+v|^2=2(|p|^2+|v|^2)$, prove the isometry. For arbitrary $w$,
the inverse is
\begin{align}
 p(x)&=\frac{w(\ell+x/c)+w(\ell-x/c)}{\sqrt{2c}},&
 v(x)&=\frac{w(\ell+x/c)-w(\ell-x/c)}{\sqrt{2c}},\\
 u(x)&=c^{-1}\int_0^x p(y)\dd y.&&
\end{align}
This lies in $\mathcal H$. For domain vectors the two midpoint traces agree
because $v(0)=0$. At the outer endpoints, $p(L)=-\eta v(L)$ gives
$w(\tau)=qw(0)$, without division by $q$. Conversely those trace relations
recover the physical boundaries and Sobolev regularity. Differentiation on
each half proves $UG_\eta U^{-1}=A_q$.
\end{proof}
In particular the dissipativity identity retains the physical constant:
\begin{equation}
 \Ree\langle A_qw,w\rangle
 =\tfrac12(q^2-1)|w(0)|^2=-c\eta|v(L)|^2.
\end{equation}

\section{The complete evolution and spectrum}
\begin{theorem}[All times and all spectral points]
For $t\geq0$, set $k=\lfloor(s+t)/\tau\rfloor$ for almost every
$s\in(0,\tau)$. The strongly continuous semigroup and its exact norm are
\begin{equation}
 S_q(t)w(s)=q^kw(s+t-k\tau),\qquad
 \|S_q(t)\|=|q|^{\lfloor t/\tau\rfloor}.
 \label{eq:semigroup}
\end{equation}
At $q=0$, the coefficient means one for $k=0$ and zero for $k\geq1$.
If $q\ne0$, the entire spectrum is algebraically simple:
\begin{equation}
 \lambda_n=\frac{\log|q|+\ii(\arg q+2\pi n)}{\tau},\qquad n\in\Z.
 \label{eq:spectrum}
\end{equation}
The functions $\tau^{-1/2}\e^{\lambda_ns}$ form a complete Riesz basis.
If $q=0$, the spectrum is empty and the sharp extinction time is $\tau$.
\end{theorem}
\begin{proof}
Characteristics give~\eqref{eq:semigroup}. Floor addition proves the
semigroup law; translation continuity on a dense set, followed by the
contraction estimate, proves strong continuity. For $t=n\tau+a$ with
$0\leq a<\tau$, the two translated input intervals have disjoint supports
and weights $q^n,q^{n+1}$. The first has positive length $\tau-a$.
Integration proves the norm bound, and data supported on that first input
interval attain it. This includes the exact endpoints $t=n\tau$.

Solving $zw-w'=f$ gives the kernel
\begin{equation}
 R(z,A_q)(s,r)=\e^{z(s-r)}
 \left[\frac{\e^{z\tau}}{\e^{z\tau}-q}-{\bf1}_{r<s}\right],
 \qquad \e^{z\tau}\ne q.
 \label{eq:green}
\end{equation}
It is a bounded kernel mapping into~\eqref{eq:domain}; direct substitution
verifies both inverse identities. At each excluded point the exponential
is an eigenvector, so no residual or continuous spectral points are omitted.
For $\alpha=(\log|q|+\ii\arg q)/\tau$, multiplication by $\e^{\alpha s}$
conjugates the periodic derivative plus $\alpha$ to $A_q$.
The Fourier basis then proves the Riesz and algebraic-simplicity statements.
At $q=0$,~\eqref{eq:green} exists everywhere and reduces to
\begin{equation}
 R(z,A_0)f(s)=\int_s^\tau\e^{z(s-r)}f(r)\dd r,
 \qquad z\in\C.
 \label{eq:volterra}
\end{equation}
Its locally uniform exponential series is operator-norm entire.
Equation~\eqref{eq:semigroup} is norm one for $t<\tau$ and zero for
$t\geq\tau$, proving sharp extinction independently of spectral language.
\end{proof}
The condition number of the specified multiplication similarity is
$1/|q|$, not an asserted optimal condition number among all possible bases.
For $q=-1$ the group is unitary. Compact resolvent and empty spectrum are
compatible for an unbounded non-self-adjoint generator; the nonempty-spectrum
theorem for bounded complex operators does not apply to $A_0$.

\begin{table}[ht]
\centering
\begin{tabular}{cccc}
\toprule
$\eta$ & $q$ & $\|S_q(n\tau)\|$, $n\geq1$ & Imaginary spectral ladder\\
\midrule
$0$ & $-1$ & $1$ & $(2n+1)\pi/\tau$\\
$1/2$ & $-1/3$ & $3^{-n}$ & $(2n+1)\pi/\tau$\\
$1$ & $0$ & $0$ & No spectral points\\
$2$ & $1/3$ & $3^{-n}$ & $2n\pi/\tau$\\
\bottomrule
\end{tabular}
\caption{Exact parameter controls. Equal decay norms do not determine the
spectral ladder. The ladder index ranges over all integers, independently
of the positive time-step index in the norm column.}
\end{table}

\ifnum\CRevisionRound>0
\section{Exact transparent pseudospectra}
Write $z=x+\ii y$. Multiplication by $\e^{\ii ys}$ is unitary and
conjugates $R(x,A_0)$ to $R(z,A_0)$. Thus the resolvent norm is a function
$\rho_\tau(x)$ of the real part alone. Empty spectrum does not make that
function vanish or remain small.

\begin{theorem}[The three branches of the exact norm]
For $x\tau>-1$, let $\theta\in(0,\pi)$ be the unique solution of
$x\tau=-\theta\cot\theta$. For $x\tau<-1$, let $h>0$ be the unique
solution of $x\tau=-h\coth h$. Then
\begin{equation}
 \rho_\tau(x)=
 \begin{cases}
 \tau\sin\theta/\theta,&x\tau>-1,\\
 \tau,&x\tau=-1,\\
 \tau\sinh h/h,&x\tau<-1.
 \end{cases}
 \label{eq:norm}
\end{equation}
It is continuous and strictly decreasing from infinity to zero.
With the strict convention $\|R\|>\varepsilon^{-1}$, every pseudospectrum is
\begin{equation}
 \sigma_\varepsilon(A_0)=\{z\in\C:\Ree z<a_\varepsilon\},\qquad
 \rho_\tau(a_\varepsilon)=\varepsilon^{-1}.
\end{equation}
\end{theorem}
\begin{proof}
The reciprocal squared norm is the infimum of
$\|xu-u'\|^2/\|u\|^2$ over nonzero $u\in H^1$ with $u(\tau)=0$.
Expansion of the numerator gives
\begin{equation}
 \int_0^\tau(|u'|^2+x^2|u|^2)\dd s+x|u(0)|^2.
\end{equation}
Compact interval embedding gives a minimizer, and variation gives
\begin{equation}
 -u''+x^2u=\mu u,\qquad u(\tau)=0,\quad u'(0)=xu(0).
 \label{eq:sl}
\end{equation}
The positive functions $\sin(\theta(1-s/\tau))$,
$\tau-s$ and $\sinh(h(1-s/\tau))$ solve the three regimes.
Their eigenvalues are respectively
$\theta^2/(\tau^2\sin^2\theta)$, $1/\tau^2$ and
$h^2/(\tau^2\sinh^2h)$.

These are the lowest roots, not arbitrary solutions of the scalar boundary
equation. A minimizing eigenfunction can be taken nonnegative. ODE uniqueness
makes it strictly positive in the interior; orthogonality excludes a distinct
eigenvalue with another positive eigenfunction. Uniqueness of the branch
parameters follows from the positive derivatives of
$-\theta\cot\theta$ and $h\coth h$.
Finally, $\sin\theta/\theta$ strictly decreases and $\sinh h/h$ strictly
increases. Their limits at zero are one. Together with the endpoint limits
of the parameter equations, this proves strict monotonicity, continuity and
the entire half-plane assertion.
\end{proof}
In particular $\rho_\tau(0)=2\tau/\pi$. The critical branch point is
$x=-1/\tau$, not zero; at that point the norm equals $\tau$.
The theorem is restricted to the transparent face. It does not assert the
same half-plane shape for a nonzero reflection coefficient.
\fi

\ifnum\CRevisionRound>1
\section{Operator ideals and the determinant boundary}
\begin{proposition}[Compact is not trace class]
For every $z\in\C$, the transparent resolvent is Hilbert--Schmidt but not
trace class, is quasinilpotent, and satisfies
\begin{equation}
 \det{}_2(I-wR(z,A_0))=1,\qquad w\in\C.
 \label{eq:det2}
\end{equation}
The semigroup is noncompact for $0\leq t<\tau$ and zero thereafter.
For $q\ne0$ it is noncompact at every time.
\end{proposition}
\begin{proof}
At $x=0$ the complete separated problem~\eqref{eq:sl} gives
\begin{equation}
 s_n(R(0,A_0))=\frac{\tau}{\pi(n+1/2)},\qquad n=0,1,2,\ldots.
 \label{eq:singular}
\end{equation}
The associated half-integer Fourier bases are complete by reflection.
The singular values are square summable and not summable.
For general $z$, $R(z,A_0)=M_zR(0,A_0)M_z^{-1}$, where $M_z$ is bounded
invertible multiplication by $\e^{zs}$. The two-sided ideal property
preserves membership, not individual singular values. Direct integration gives
\begin{equation}
 \|R(x+\ii y,A_0)\|_{\rm HS}^2=
 \begin{cases}
 \displaystyle\frac{\tau}{2x}-\frac{1-\e^{-2x\tau}}{4x^2},&x\ne0,\\
 \tau^2/2,&x=0.
 \end{cases}
\end{equation}
The $n$th power of the Volterra inverse has kernel
$\e^{z(s-r)}(r-s)^{n-1}/(n-1)!$ for $r>s$.
Young's inequality implies
$\|R(z,A_0)^n\|\leq\e^{|x|\tau}\tau^n/n!$, hence spectral radius zero.
The order-two determinant is explicitly the canonical product
$\prod_j(1-w\lambda_j)\e^{w\lambda_j}$ over nonzero compact-operator
eigenvalues with algebraic multiplicity. There are no such eigenvalues here,
so that product is one.

For $t<\tau$, the evolution is isometric on inputs supported on $(t,\tau)$.
An orthonormal sequence on that infinite-dimensional subspace has no
convergent image subsequence, proving noncompactness. For $q\ne0$, the
wrapped translation has a bounded inverse at every fixed time; compactness
would make the identity compact on an infinite-dimensional space.
\end{proof}
Thus there is no ordinary trace of the semigroup before transparent
extinction and no ordinary trace-class Fredholm determinant of its
transparent resolvent. After extinction the semigroup is zero and has trace
zero. The scalar characteristic $1-q\e^{-z\tau}$ records the nontransparent
spectral equation but is not thereby an ordinary operator determinant.
Equation~\eqref{eq:det2} names its regularization; it does not forbid all
other regularized constructions.

\paragraph{Parameter and reversal controls.}
For positive impedance, $q(1/\eta)=-q(\eta)$. This preserves the norm and
shifts the spectral ladder by $\ii\pi/\tau$ modulo $2\pi\ii/\tau$.
As $\eta\to1$, the real parts of every nontransparent eigenvalue tend to
minus infinity, while the similarity condition number diverges. Nevertheless
the norm stays one at every fixed $t<\tau$. The physical reversal
$(u,v)\mapsto(u,-v)$ maps the impedance domain for $\eta$ to that for
$-\eta$, and therefore does not preserve positive damping. The conservative
endpoint $\eta=0$ has the usual physical time reversal. These statements
concern the physical involution, not arbitrary abstract antiunitaries.
\fi

\section{Executable evidence and scope}
The canonical ledger uses seven impedances, three round-trip times and
seven time ratios. Four spatial ratios give 588 exact transport rows;
seven boundary rows check the reflection sign and dissipation. A separate
checker uses repeated boundary crossings instead of the producer's floor
formula. It checks 126 finite spectral rows by the exponential boundary
equation, and 21 Green rows by direct kernel integration instead of the
producer's polynomial inhomogeneous solution.
\ifnum\CRevisionRound>0
The 27 transparent norm rows include all three regimes and the zero-real-part
case. A separate 100-digit lane evaluates 27 Rayleigh quotients,
81 Volterra action checks and 81 complex-gauge action checks. Eleven
symbolic identities check the characteristic, boundary and differential
relations. Finite regression of 12 half-integer singular modes supports the
normalization; the complete sequence follows from the Fourier argument.
\fi
The stored decimal values have 60 significant digits from 100-digit
arithmetic. They are not interval certificates. Canonical reproduction,
two-directory replay and repaired-hash mutations test the artifact contract;
they cannot establish the all-time and full-spectrum theorem by enumeration.
Unknown fields and Boolean--integer substitutions are rejected by typed
checks. The frozen evaluation is parsed strictly and locked before a
release-write operation. Compiler logs, font checks and PDF hashes are
separate engineering evidence, not mathematical novelty.

\begin{samepage}
The conservative Route-A tuple is four failures followed by a formal
operator hint; overall the candidate is rejected for the target route.
The native generator supplies no rational-prime carrier, arithmetic
primitive correspondence, target divisor, target functional equation or
Hilbert--P\'olya identification. No Euler factors, root number, automorphy
or target local data are claimed. Route B remains disabled.
\par\noindent\texttt{NO\_BAD\_EULER\_OR\_ROOT\_NUMBER}.
\end{samepage}

\paragraph{Literature and review boundaries.}
The author's website PDF initially failed to download. Retrieval from the
same author's GitHub repository succeeded, and the relevant extracted text
was read. The optional structural preflight was unavailable because its
PDF-parser dependency was absent; no certified local-page-anchor claim is
made. Equations (14)--(19), (25) and Theorem~1 of~\cite{driscoll1996}
support the ownership statement.
\ifnum\CRevisionRound=0
The three-branch norm is established in the separate complete proof package;
it is not attributed to a formula absent from the cited theorem.
\else
The precise three-branch norm is derived above rather than attributed to a
formula not present in that theorem.
\fi
An internal team agent independently checked the complete analytic proof.
The authoring agent supplied the implementation and manuscript. This is
AI-assisted internal review, not external peer review or venue acceptance.
Variable coefficients, internal damping and infinite impedance are outside
the frozen theorem. They would require new domains and proofs.

\section{Conclusion}
\ifcase\CRevisionRound
\textbf{Round zero: the physical domain and complete spectrum.}
The same exact unfolding determines all spectral points and all evolution
times. At transparency the generator has no spectrum, yet the norm remains
one until the sharp physical round-trip time.
\or
\textbf{Round one: transparent resolvent and exact pseudospectra.}
The lowest separated eigenfunction selects three explicit resolvent-norm
branches. Their strict monotonicity gives every transparent pseudospectrum,
with no reliance on a finite matrix's apparent eigenvalues.
\else
\textbf{Round two: operator ideals and the determinant boundary.}
Finite extinction, empty spectrum and an entire resolvent coexist in a
single natural PDE source. Its resolvent is compact but not trace class,
and its explicitly regularized determinant is trivial. These distinctions
close the source theorem while leaving the target arithmetic route rejected.
\fi

\begin{thebibliography}{1}
\raggedright
\bibitem{driscoll1996}
T. A. Driscoll and L. N. Trefethen,
Pseudospectra for the wave equation with an absorbing boundary,
\emph{Journal of Computational and Applied Mathematics} \textbf{69}
(1996), 125--142.
\href{https://doi.org/10.1016/0377-0427(95)00021-6}{doi:10.1016/0377-0427(95)00021-6}.
\end{thebibliography}
\end{document}
